\section{Introduction}
%% Storyline (unified replacement): a text-only teacher says what to
%% do, not what to replace; imitation replaces wholesale;
%% return-targeted distillation gives every sampled student behavior
%% its own share to hand to the demonstration and chooses the shares
%% jointly by the learning return; evidence is priced by the gain it
%% adds to the same teaching problem. House style: no em dashes, no colons, no prose
%% parentheticals, one step of logic per sentence.

Language-model agents are increasingly deployed for narrow, recurring
jobs, such as operating a fixed set of applications, handling customer
requests over database-backed tools, or automating one workflow
\citep{yao2023react, trivedi2024appworld, yao2024taubench,
shridhar2021alfworld}. Serving a frontier model for a single narrow
job is rarely economical, so practice converges on specializing a
small student for that job. The deployer typically has a few dozen
example queries from the intended workload, a task-level check of
whether an attempt succeeded, and metered query access to a strong
commercial teacher whose weights and logits are out of reach and whose
every emitted token is billed. This paper asks how a limited teacher
budget can be converted into the largest possible gain in task success
for the small student.

The standard answer is imitation. One purchases teacher demonstrations
and trains the student to reproduce them, the recipe behind
sequence-level knowledge distillation and its agentic refinements
\citep{kim2016sequence, mukherjee2023orca, liu2025structured,
kang2025distilling}. A demonstration from a text-only teacher,
however, answers only half of the question the student faces. It says
what the teacher would do at a state. It does not say which of the
existing behaviors of the student at that state should be replaced and
which should be kept, and maximum-likelihood training replaces all of
them at once, since its expected update descends the divergence between
the teacher and student behavior distributions as a whole. On agent
benchmarks the student already solves part of the task, and the part
it solves is exactly what wholesale replacement puts at risk.
Preference-based and on-policy variants change what is compared, but
they still let the teacher text decide both how much to move and what
to move away from.

We make the two decisions one. Return-targeted distillation (RTD)
installs the teacher text through a source-to-teacher replacement. At
every supervised state the student first exhibits its own behaviors
by sampling from a frozen snapshot, and each sampled behavior receives
its own share to hand to the teacher demonstration; the teacher dose
and the choice of what gives way are read off the shares rather than
decided separately. The shares are not read off the teacher text. They
are chosen jointly for a whole training step by the learning return
of the student, measured after a virtual update on the same batch,
through one convex problem in which every share is a concrete
direction, the replacement of one student behavior by one
demonstration, and an interaction term charges directions that push
the update the same way. The estimator that installs the chosen
target keeps the expectation of the target gradient even though the
shares depend on the very samples. The same problem prices
acquisition. A candidate purchase is worth the gain it adds to what
the evidence already owned can teach, including the teaching that
only the new demonstration makes possible, and purchases are executed
through a hard ledger that bills failed completions at their recorded
cost under a budget that is a cap rather than an amount to be spent.

This paper makes three contributions. The first is a problem
statement. Budgeted few-shot agent specialization asks, given a
few-shot description of an agent task and a budget on a per-token
billed, text-only black-box teacher, for a small model that maximizes
task performance; the budget is the single constraint and held-out
task performance the single criterion. The second is a mechanism. We
show that the per-source replacement target is a legal distribution
whose gradient estimator keeps its expectation under sample-dependent
shares, that the resulting update is affine in the shares with the
natural degeneracies, and that source selection is a strictly larger
family than a per-state teacher dose; the practical selection is a
convex local problem with an explicit interaction term and an
uncertainty penalty, and the same problem defines the value of
evidence. The
third is a protocol with results. Under a common teacher budget cap per
benchmark, with the agent harness fixed and only the student backbone
changing, we compare RTD with implementation-faithful distillation
baselines, and we isolate the mechanism with paired controls that hold
the purchased data and the source samples fixed and change only the
replacement decision. \todo{finalize headline numbers
at freeze.}
